\section{EU AI Act}

Algunas de las cosas que menciona el EU AI Act \cite{EU2024_1689} sobre la apertura de los LLMs son:


"Software and data, including models, released under a free and open-source licence that allows them to be openly shared and where users can freely access, use, modify and redistribute them or modified versions thereof, can contribute to research and innovation in the market and can provide significant growth opportunities for the Union economy. General-purpose AI models released under free and open-source licences should be considered to ensure high levels of transparency and openness if their parameters, including the weights, the information on the model architecture, and the information on model usage are made publicly available. The licence should be considered to be free and open-source also when it allows users to run, copy, distribute, study, change and improve software and data, including models under the condition that the original provider of the model is credited, the identical or comparable terms of distribution are respected." 

"Free and open-source AI components covers the software and data, including models and general-purpose AI models, tools, services or processes of an AI system. Free and open-source AI components can be provided through different channels, including their development on open repositories. For the purposes of this Regulation, AI components that are provided against a price or otherwise monetised, including through the provision of technical support or other services, including through a software platform, related to the AI component, or the use of personal data for reasons other than exclusively for improving the security, compatibility or interoperability of the software, with the exception of transactions between microenterprises, should not benefit from the exceptions provided to free and open-source AI components. The fact of making AI components available through open repositories should not, in itself, constitute a monetisation." 

"The providers of general-purpose AI models that are released under a free and open-source licence, and whose parameters, including the weights, the information on the model architecture, and the information on model usage, are made publicly available should be subject to exceptions as regards the transparency-related requirements imposed on general-purpose AI models, unless they can be considered to present a systemic risk, in which case the circumstance that the model is transparent and accompanied by an open-source license should not be considered to be a sufficient reason to exclude compliance with the obligations under this Regulation. In any case, given that the release of general-purpose AI models under free and open-source licence does not necessarily reveal substantial information on the data set used for the training or fine-tuning of the model and on how compliance of copyright law was thereby ensured, the exception provided for general-purpose AI models from compliance with the transparency-related requirements should not concern the obligation to produce a summary about the content used for model training and the obligation to put in place a policy to comply with Union copyright law, in particular to identify and comply with the reservation of rights pursuant to Article 4(3) of Directive (EU) 2019/790 of the European Parliament and of the Council (40).


Esta va para las menciones de copy right:

"General-purpose AI models, in particular large generative AI models, capable of generating text, images, and other content, present unique innovation opportunities but also challenges to artists, authors, and other creators and the way their creative content is created, distributed, used and consumed. The development and training of such models require access to vast amounts of text, images, videos and other data. Text and data mining techniques may be used extensively in this context for the retrieval and analysis of such content, which may be protected by copyright and related rights. Any use of copyright protected content requires the authorisation of the rightsholder concerned unless relevant copyright exceptions and limitations apply. Directive (EU) 2019/790 introduced exceptions and limitations allowing reproductions and extractions of works or other subject matter, for the purpose of text and data mining, under certain conditions. Under these rules, rightsholders may choose to reserve their rights over their works or other subject matter to prevent text and data mining, unless this is done for the purposes of scientific research. Where the rights to opt out has been expressly reserved in an appropriate manner, providers of general-purpose AI models need to obtain an authorisation from rightsholders if they want to carry out text and data mining over such works." 

Esto para la transparencia de los datos de entrenamiento:

"In order to increase transparency on the data that is used in the pre-training and training of general-purpose AI models, including text and data protected by copyright law, it is adequate that providers of such models draw up and make publicly available a sufficiently detailed summary of the content used for training the general-purpose AI model. While taking into due account the need to protect trade secrets and confidential business information, this summary should be generally comprehensive in its scope instead of technically detailed to facilitate parties with legitimate interests, including copyright holders, to exercise and enforce their rights under Union law, for example by listing the main data collections or sets that went into training the model, such as large private or public databases or data archives, and by providing a narrative explanation about other data sources used. It is appropriate for the AI Office to provide a template for the summary, which should be simple, effective, and allow the provider to provide the required summary in narrative form."